\indent “ಕುಮಾರು ಓಡುತ್ತದೆ” ಎಂಬ ವಾಕ್ಯದಲ್ಲಿ ಕರ್ಮಪದವಿಲ್ಲ. ಎಂಬ ಧಾತುವಿಗೆ ಕರ್ಮಪದದ ಅಗತ್ಯ ಬೀಳುವುದಿಲ್ಲ... ಹೀಗೆ, ಕರ್ಮ ಅಗತ್ಯ ಇರದ ಧಾತುವೇ ಅಕರ್ಮಕ ಧಾತು\par

\indent “ನಾಯಿಯು ರೈತನು ತಿನ್ನುತ್ತದೆ”--ಇದು ಕರ್ಮಪದ ಹೊಂದಿರುವ ವಾಕ್ಯ. “ತಿನ್ನು” ಎಂಬ ಧಾತುವಿಗೆ ಕರ್ಮಪದದ ಅಗತ್ಯವ ಹೀಗೆ ಕರ್ಮಪದದ ಅಗತ್ಯವಿರುವ ಧಾತುವೇ ಸಕರ್ಮಕ ಧಾತು.\par

\indent ಉದಾಹರಣೆಗೆ ಬೀಳು, ಏಕ ಕುಣಿ, 'ಹೋಗು, ಬರು, ಅಕರ್ಮಕ ಧಾತುಗಳು ; ಕಾಣು, ಮಾಡು, ನೋಡು, 'ಹಾಡು 8 ಕೊಡು, ಸಕರ್ಮಕ ಧಾತುಗಳು. (1 ೫1 ಗ\par

\kbvexercise{ಅಭ್ಯಾಸ ಪತ್ರಿಕೆ ೫}

\vspace{2.5pt plus 3pt minus 0.5pt}\noindent\hspace*{1.2em}\kbvnum{1.}\hspace{0.55em}ಕ್ರಿಯಾಪದ ಎಂದರೇನು ? ಉದಾಹರಣೆ ಸಹಿತ ತಿಳಿಸಿ.\par

\vspace{2.5pt plus 3pt minus 0.5pt}\noindent\hspace*{1.2em}\kbvnum{2.}\hspace{0.55em}ಕುಣಿಯುತ್ತದೆ, ಮಾಡುತ್ತದೆ, ಬರೆಯುತ್ತಾಳೆ- ಈ ಕ್ರಿಯಾ ಗಳಲ್ಲಿ ಪ್ರಕೃತಿ, ಪ್ರತ್ಯಯಗಳನ್ನು ಬಿಡಿಸಿ ಬರೆಯಿರಿ.\par

\vspace{2.5pt plus 3pt minus 0.5pt}\noindent\hspace*{1.2em}\kbvnum{3.}\hspace{0.55em}ಧಾತು ಎಂದರೇನು ? ಎರಡು ಉದಾಹರಣೆಗಳನ್ನು ಕೊಡಿ.\par

\vspace{2.5pt plus 3pt minus 0.5pt}\noindent\hspace*{1.2em}\kbvnum{4.}\hspace{0.55em}ಆಖ್ಯಾತ ಪ್ರತ್ಯಯಗಳಿಗೆ ಎರಡು ಉದಾಹರಣೆಗಳನ್ನು ಕೊಡಿ\par

\vspace{2.5pt plus 3pt minus 0.5pt}\noindent\hspace*{1.2em}\kbvnum{5.}\hspace{0.55em}ಅಕರ್ಮಕ ಧಾತುವಿಗೂ ಸಕರ್ಮಕ ಧಾತುವಿಗೂ' ಇಃ ವ್ಯತ್ಯಾಸವೇನು ? ಉದಾಹರಣೆಗಳ ಸಹಿತ ತಿಳಿಸಿ.\par

\vspace{2.5pt plus 3pt minus 0.5pt}\noindent\hspace*{1.2em}\kbvnum{6.}\hspace{0.55em}ಕರ್ಮಪದವನ್ನು ಅಪೇಕ್ಷಿಸುವ ಎರಡು ಧಾತುಗಳನ್ನು ತಿಳಿಸಿ,\par

\vspace{2.5pt plus 3pt minus 0.5pt}\noindent\hspace*{1.2em}\kbvnum{7.}\hspace{0.55em}ಕರ್ಮಪದವನ್ನು ಅಪೇಕ್ಷಿಸದ ಎರಡು ಧಾತುಗಳನ್ನು ತಿಳಿಸಿ.\par

\vspace{2.5pt plus 3pt minus 0.5pt}\noindent\hspace*{1.2em}\kbvnum{8.}\hspace{0.55em}ಪ್ರತ್ಯಯಾಂತ ಧಾತುವಿಗೆ ಒಂದು ಉದಾಹರಣೆಯನ್ನು ಕೊಡಿ\par

\vspace{2.5pt plus 3pt minus 0.5pt}\noindent\hspace*{1.2em}\kbvnum{9.}\hspace{0.55em}“ದುಃಖಿಸುತ್ತಾನೆ” ಈ ಕ್ರಿಯಾಪದದಲ್ಲಿ ಮೂರು ಪ್ರತ್ಯಯ ಸೇರಿವೆ ; ಅವುಗಳನ್ನು ಬಿಡಿಸಿ ಹೇಳಿ.\par

\vspace{2.5pt plus 3pt minus 0.5pt}\noindent\hspace*{1.2em}\kbvnum{10.}\hspace{0.55em}ಪ್ರಕೃತಿ, ಪ್ರತ್ಯಯಗಳನ್ನು ಸೇರಿಸಿ ಪದಗಳನ್ನು ಬರೆಯಿರಿ.\par

\vspace{2.5pt plus 3pt minus 0.5pt}\noindent\hspace*{1.2em}\kbvnum{(೧)}\hspace{0.55em}ಕಲಿ + ಉತ್ತ + ಆರೆ (೨) ಮಾಡು + ದ + ಅನು\par

\vspace{2.5pt plus 3pt minus 0.5pt}\noindent\hspace*{1.2em}\kbvnum{(೩)}\hspace{0.55em}ನೋಡು + ದ + ಅಳು (೪) ಮಲಗು + ಇಸು\par
